\documentclass[11pt,a4paper]{article}
\usepackage[UTF8]{ctex}
\usepackage{amsmath,amssymb}
\usepackage{geometry}
\usepackage{hyperref}
\geometry{margin=1in}
\title{PPO 与 AMP 复习笔记}
\author{}
\date{}

\begin{document}
\maketitle

这份笔记写给“已经学过，但一段时间没看细节”的自己。目标不是替代教科书，而是帮你在做项目、看日志、调参数时，快速想起：

\begin{itemize}
\item PPO 在优化什么
\item AMP 比 PPO 多了什么
\item 为什么 AMP 更容易只会站、不愿走
\item 日志里那些量到底该怎么理解
\end{itemize}

本文尽量结合你当前这个项目：
\begin{itemize}
\item \texttt{/root/Desktop/Bipedal\_AMP}
\item Isaac Lab + forked rsl\_rl + AMP
\end{itemize}

\section{先记住一句话}
\subsection{PPO}
PPO 是一种“限制每次策略更新不要太猛”的策略梯度方法。

\subsection{AMP}
AMP = PPO 的任务学习 + 一个“动作像不像参考数据”的判别器奖励。

所以：
\begin{itemize}
\item PPO 只学“完成任务”
\item AMP 同时学“完成任务”和“动作风格像参考数据”
\end{itemize}

\section{PPO 到底在学什么}
在强化学习里，我们想让策略 $\pi_\theta(a|s)$ 最大化长期回报：
\[
J(\theta)=\mathbb{E}_{\tau \sim \pi_\theta}\left[\sum_t \gamma^t r_t\right]
\]

其中：
\begin{itemize}
\item $s$：状态
\item $a$：动作
\item $r_t$：奖励
\item $\gamma$：折扣因子
\end{itemize}

但是直接暴力更新策略容易不稳定，所以 PPO 的核心思想是：

\begin{quote}
允许策略变好，但不允许一次改太多。
\end{quote}

\section{PPO 的核心式子}
PPO 最经典的是 clipped objective：
\[
L^{\text{CLIP}}(\theta)
=
\mathbb{E}\left[
\min\left(
r_t(\theta)\hat A_t,\;
\text{clip}(r_t(\theta),1-\epsilon,1+\epsilon)\hat A_t
\right)
\right]
\]

其中：
\[
r_t(\theta)=
\frac{\pi_\theta(a_t|s_t)}
{\pi_{\theta_{\text{old}}}(a_t|s_t)}
\]

$\hat A_t$ 是 advantage，表示这一步动作到底比“平均水平”好多少。

\subsection{直观理解}
\begin{itemize}
\item 如果新策略比旧策略只变一点点，那就正常更新
\item 如果新策略偏得太狠，就把更新裁掉
\end{itemize}

这就是 PPO 的“保守更新”。

\section{PPO 里 value function 是干什么的}
PPO 一般是 actor-critic：
\begin{itemize}
\item actor 负责输出动作分布
\item critic 负责估计状态价值 $V(s)$
\end{itemize}

优势函数常写成：
\[
A_t = Q(s_t,a_t)-V(s_t)
\]

意思是：
\begin{itemize}
\item 如果这个动作比当前状态的平均预期更好，$A_t>0$
\item 如果更差，$A_t<0$
\end{itemize}

实践里 PPO 常用 GAE：
\[
\hat A_t^{\text{GAE}}=\sum_{l=0}^{\infty}(\gamma\lambda)^l\delta_{t+l}
\]

其中 TD 误差：
\[
\delta_t=r_t+\gamma V(s_{t+1})-V(s_t)
\]

\subsection{记忆版理解}
\begin{itemize}
\item critic 不是为了控制机器人
\item 它是为了帮助 actor 更稳定地知道“这一步到底值不值”
\end{itemize}

\section{PPO 的总损失通常长什么样}
一个常见的 PPO 总损失可以写成：
\[
L_{\text{PPO}} = L_{\text{policy}} + c_v L_{\text{value}} - c_e H(\pi)
\]

其中：
\begin{itemize}
\item policy loss：希望策略变好
\item value loss：让 critic 估值更准
\item entropy：鼓励探索，不要过早变得太确定
\end{itemize}

\subsection{熵项的直觉}
\begin{itemize}
\item 熵太低：策略太早“自信”，容易陷入局部最优
\item 熵太高：策略一直乱试，学不稳
\end{itemize}

\section{训练日志里 PPO 指标怎么理解}
\subsection*{Mean reward}
\begin{itemize}
\item 越高通常越好
\item 但必须结合 reward 构成看
\item 单看一个总 reward 容易被误导
\end{itemize}

\subsection*{Mean episode length}
\begin{itemize}
\item 变长通常说明更不容易摔
\item 但不代表任务就学好了
\item 很可能只是“会站住了”
\end{itemize}

\subsection*{value loss}
\begin{itemize}
\item 反映 critic 学得好不好
\item 不是越小越绝对好，但太大通常说明估值不稳
\end{itemize}

\subsection*{surrogate loss}
\begin{itemize}
\item PPO 的策略目标
\item 绝对值小不一定坏
\item 要结合 reward 趋势看
\end{itemize}

\subsection*{entropy loss / mean action noise std}
\begin{itemize}
\item 反映探索程度
\item 太低：容易卡住
\item 太高：动作可能太乱
\end{itemize}

\section{PPO 为什么常常比 AMP 更容易收敛}
因为 PPO 只需要解决：
\begin{quote}
怎么做能拿更高任务奖励？
\end{quote}

而 AMP 需要同时解决：
\begin{quote}
怎么做能拿更高任务奖励？\\
怎么做还得像参考动作？
\end{quote}

这让优化目标变得更复杂，也更容易互相冲突。

\section{AMP 到底比 PPO 多了什么}
AMP 的关键额外部分是：
\begin{enumerate}
\item 判别器 $D$
\item 参考动作数据 demo
\item style reward
\end{enumerate}

\subsection{核心思想}
判别器学习区分：
\begin{itemize}
\item demo 里的参考动作片段
\item 当前策略产生的动作片段
\end{itemize}

如果判别器觉得策略动作“像 demo”，就给策略更高的 style reward。

所以 AMP 不是直接监督学习，而是：
\begin{quote}
用一个对抗式判别器，把“像不像参考动作”转成奖励。
\end{quote}

\section{AMP 的直观结构}
可以把 AMP 理解成：
\[
r_t^{\text{total}}
=
(1-\alpha)r_t^{\text{task}}
+
\alpha r_t^{\text{style}}
\]

这里：
\begin{itemize}
\item $r_t^{\text{task}}$：任务奖励，比如速度跟踪、别摔、动作平滑
\item $r_t^{\text{style}}$：风格奖励，由判别器给
\item $\alpha$：风格和任务的混合权重
\end{itemize}

在你这个项目里，对应的是 \texttt{task\_style\_lerp} 这类配置。

\subsection{这就是为什么 AMP 常见现象是}
\begin{itemize}
\item 先学会站稳
\item 动作“有点像”
\item 但并不一定马上会走
\end{itemize}

因为 style 本身就可能奖励“看起来比较像站姿/基础节律”的解。

\section{判别器在 AMP 里学什么}
判别器不是看单帧就够，而通常看一个短时间窗口：
\[
o_t^{\text{disc}} = [x_t, x_{t-1}, \dots, x_{t-k}]
\]

也就是说，判别器看的是“短时间动态”，不是只看一个姿势。

所以它能区分：
\begin{itemize}
\item 一个姿势像不像
\item 动作变化节律像不像
\end{itemize}

在你当前项目里，\texttt{disc} 和 \texttt{disc\_demo} 这两组观测就是干这个的。

\section{AMP 判别器的损失在干什么}
虽然不同实现细节略有区别，但核心目标都是：
\begin{itemize}
\item 让 demo 被判成“真”
\item 让 policy 产生的 motion 被判成“假”
\end{itemize}

如果更抽象地写，判别器在最大化：
\[
\log D(x_{\text{demo}}) + \log(1-D(x_{\text{policy}}))
\]

实践里还常带：
\begin{itemize}
\item gradient penalty
\item weight decay
\end{itemize}

就是为了让判别器更稳，不要太容易爆。

\section{为什么 AMP 更容易“只会站、不愿走”}
这是做双足 AMP 时最常见的现象之一。

原因通常是这几个一起造成的：
\begin{enumerate}
\item 任务奖励和风格奖励冲突
\item 站立是低风险解
\item 终止条件太宽松
\item 参考数据质量有问题
\end{enumerate}

例如：
\begin{itemize}
\item 任务想让它往前走
\item style 数据却更像原地踏步或低速动作
\end{itemize}

\section{为什么“会站但不走”不能简单理解为训练正常}
如果一个 AMP 策略：
\begin{itemize}
\item episode length 很高
\item time\_out 很多
\item 但 track\_lin\_vel\_xy\_exp 很低
\item error\_vel\_xy 还很大
\end{itemize}

那更像是在学：
\begin{quote}
怎样安全地活到 episode 结束
\end{quote}

而不是：
\begin{quote}
怎样按照命令稳定行走
\end{quote}

这是你看日志时特别要警惕的局部最优。

\section{什么时候更像是“没训够”，什么时候更像是“卡住了”}
\subsection*{更像没训够}
\begin{itemize}
\item track\_lin\_vel\_xy\_exp 持续在涨
\item error\_vel\_xy 在慢慢降
\item style 也在涨
\item episode length 增长同时任务指标也改善
\end{itemize}

\subsection*{更像卡住}
\begin{itemize}
\item episode length 很高
\item time\_out 很高
\item 但速度跟踪长期不上升
\item 机器人主要就是稳站或小幅乱挪
\end{itemize}

\section{AMP 为什么这么吃数据}
因为 AMP 学的是一个更复杂的问题：
\begin{itemize}
\item 既要满足控制目标
\item 又要满足动作先验
\item 还要让判别器和策略互相博弈
\end{itemize}

所以它通常比纯 PPO 更吃：
\begin{itemize}
\item 总采样量
\item motion data 质量
\item 训练稳定性
\end{itemize}

\section{你当前项目里 PPO 和 AMP 的结构差异}
在这个项目里，AMP 比 PPO 多出来的核心是：
\begin{itemize}
\item disc 观测
\item disc\_demo 观测
\item motion\_data\_manager
\item animation\_manager
\item amp\_discriminator
\item PPOAMP 算法
\end{itemize}

相关位置可以看：
\begin{itemize}
\item \href{/root/Desktop/Bipedal_AMP/source/legged_lab/legged_lab/tasks/locomotion/amp/amp_env_cfg.py}{amp\_env\_cfg.py}
\item \href{/root/Desktop/Bipedal_AMP/source/legged_lab/legged_lab/rsl_rl/amp_cfg.py}{amp\_cfg.py}
\item \href{/root/Desktop/Bipedal_AMP/source/legged_lab/legged_lab/rsl_rl/rl_cfg.py}{rl\_cfg.py}
\end{itemize}

\section{什么时候要怀疑 motion data}
如果出现这些现象，优先怀疑 motion data：
\begin{itemize}
\item style reward 有值，但动作看起来很怪
\item joint\_pos\_penalty 特别大
\item 某些关节长期贴边
\item 训练容易数值不稳定
\item 机器人能站住，但步态很别扭
\end{itemize}

这通常意味着：
\begin{itemize}
\item retarget 坐标系没对好
\item joint 正方向没对好
\item 默认姿态和数据零位不一致
\end{itemize}

\section{标准差、探索和数值稳定}
策略常用高斯分布采样动作：
\[
a_t \sim \mathcal{N}(\mu_\theta(s_t), \sigma_\theta(s_t))
\]

如果 std 参数化不好，就可能出问题。

\subsection*{直接学 std}
风险：
\begin{itemize}
\item 数值可能被推到负数
\item 不稳定
\end{itemize}

\subsection*{学 log\_std}
更常见更稳：
\[
\sigma = \exp(\log \sigma)
\]

这样天然保证：
\[
\sigma > 0
\]

所以你之前碰到的：
\begin{quote}
normal expects all elements of std >= 0.0
\end{quote}

本质上就是策略分布数值出问题了。这未必说明任务一定学坏了，但一定说明训练稳定性出问题了。

\section{训练日志里 AMP 相关项怎么读}
\subsection*{Mean AMP total reward}
\begin{itemize}
\item AMP 风格侧总奖励
\item 不是整个训练唯一目标
\end{itemize}

\subsection*{Mean AMP style reward}
\begin{itemize}
\item 越高通常表示判别器更认可当前动作像 demo
\item 但要警惕“像得不一定对任务有帮助”
\end{itemize}

\subsection*{amp/disc\_loss}
\begin{itemize}
\item 判别器分类相关损失
\item 不是越小越绝对好，要看是否稳定
\end{itemize}

\subsection*{amp/disc\_grad\_penalty}
\begin{itemize}
\item 判别器梯度惩罚
\item 太大通常说明判别器比较激进
\end{itemize}

\subsection*{amp/disc\_demo\_score}
demo 被判真程度。

\subsection*{amp/disc\_score}
policy motion 被判真的程度。

\section{一份实用判断模板}
看一段训练日志时，可以按这个顺序问自己：
\begin{enumerate}
\item 它更会活了吗？
\item 它更会做任务了吗？
\item 它更像参考动作了吗？
\item 它有没有学歪？
\end{enumerate}

\section{为什么 AMP 会比 PPO 更依赖 reward 设计边界}
因为 AMP 已经有了一股“动作风格梯度”，所以任务 reward 不需要像纯 PPO 那么“什么都自己教”。

但如果 task reward 设计得太强、太硬，也会和 style 冲突。

所以 AMP 的 reward 调法和 PPO 不完全一样：
\begin{itemize}
\item PPO：更多靠 task reward 直接塑形
\item AMP：task reward 更像“任务边界条件”，style 负责动作风格
\end{itemize}

\section{如果我是做双足 locomotion，该怎么理解 PPO 和 AMP 的关系}
\subsection*{PPO}
像在教机器人：
\begin{quote}
你只要想办法走起来、别摔、按命令跑就行。
\end{quote}

\subsection*{AMP}
像在教机器人：
\begin{quote}
你不仅要走起来，还要走得像这批参考动作。
\end{quote}

\section{一页纸记忆版}
\subsection*{PPO 关键词}
\begin{itemize}
\item actor-critic
\item advantage
\item clipped objective
\item entropy
\item 保守更新
\end{itemize}

\subsection*{AMP 关键词}
\begin{itemize}
\item discriminator
\item demo motion
\item style reward
\item task + style
\item 比 PPO 更吃数据，也更依赖数据质量
\end{itemize}

\section{和你当前项目最相关的提醒}
结合你这段时间的实验，最该记住的是：
\begin{enumerate}
\item 会站 不等于 会走
\item style 有值 不等于 motion data 没问题
\item episode length 很长 不等于已经收敛
\item 关节长期贴边 往往说明 retarget 或 IK config 有问题
\item AMP 如果数据不对，会比 PPO 更容易被带歪
\end{enumerate}

\section{最后一句提醒}
当你以后再觉得“理论有点糊”时，先别急着看公式，先问自己这两个问题：
\begin{enumerate}
\item 这个方法到底在优化什么？
\item 它为什么会学成我现在看到的样子？
\end{enumerate}

只要这两个问题能回答上来，很多调参和诊断都会顺很多。

\end{document}
